\section{GRASP para VRPTW}
\label{sec:grasp}

O GRASP (\emph{Greedy Randomized Adaptive Search Procedure}) alterna construção gulosa randomizada e refinamento por busca local \cite{feo1995greedy,resende2003greedy}.
No contexto deste trabalho, as duas variantes implementadas (fixa e reativa) usam o mesmo núcleo de construção e o mesmo refinamento local (VND).

\subsection{Estrutura do método implementado}

Cada iteração executa três fases:
\begin{enumerate}
  \item seleção de $\alpha_t$ (fixo ou amostrado da distribuição reativa);
  \item construção de uma solução por inserção com RCL;
  \item aplicação de VND e atualização do incumbente.
\end{enumerate}

Além disso, o incumbente inicial é o I1 puro, para manter alinhamento com o baseline da comparação.
A parada é controlada por dois limites: número máximo de iterações ($N_{\max}$) e número máximo de iterações sem melhoria ($L$).

\begin{algorithm}[htbp]
\caption{GRASP (fixo e reativo) com VND}
\label{alg:grasp}
\begin{algorithmic}[1]
\Require instância VRPTW, modo (\texttt{fixed}/\texttt{reactive}), $N_{\max}$ e $L$
\Ensure melhor solução $S^*$
\State $S^* \leftarrow I1$
\State $noImprove \leftarrow 0$
\State inicializa distribuição reativa (se modo reativo)
\For{$t=1$ até $N_{\max}$}
  \State define $\alpha_t$ (fixo ou amostrado)
  \State $S_t \leftarrow$ construção com RCL usando $\alpha_t$
  \If{$S_t$ cobre todos os clientes}
    \State $S_t \leftarrow VND(S_t)$
  \EndIf
  \If{$S_t$ melhora $S^*$}
    \State $S^* \leftarrow S_t$; $noImprove \leftarrow 0$
  \Else
    \State $noImprove \leftarrow noImprove + 1$
  \EndIf
  \If{modo reativo e $t \bmod update = 0$}
    \State atualiza probabilidades dos $\alpha$ e zera acumuladores do bloco
  \EndIf
  \If{$noImprove \ge L$}
    \State \textbf{break}
  \EndIf
\EndFor
\State \Return $S^*$
\end{algorithmic}
\end{algorithm}

\paragraph{Texto complementar ao pseudocódigo.}
As linhas 3--4 implementam o controle de estagnação: toda iteração sem melhora do incumbente incrementa \texttt{noImprove}, e a execução para quando o limite $L$ é atingido.
Isso evita gastar orçamento em repetições pouco produtivas.

Na linha 6, a construção com RCL produz uma solução completa ou parcial.
Se parcial, a iteração é contabilizada, mas não entra em refinamento (linha 7), preservando consistência da comparação.
Se completa, a linha 8 aplica o mesmo VND usado no baseline de busca local, de modo que o diferencial do GRASP fica concentrado na estratégia de diversificação da fase construtiva.

No modo reativo, a atualização de probabilidades ocorre por blocos (\texttt{reactiveUpdate}) e não em toda iteração.
Essa escolha reduz variância na estimativa de desempenho de cada $\alpha$, evitando que decisões de probabilidade sejam influenciadas por amostras muito pequenas.

\subsection{Construção com RCL}

A construção usa inserção factível por custo incremental:
\[
 c_1(i,u,j)=d_{iu}+d_{uj}-\mu d_{ij}.
\]

Em cada passo, calcula-se $c_{\min}$ e $c_{\max}$ entre inserções factíveis e define-se:
\[
 \text{RCL}=\{(i,u,j): c_1(i,u,j) \le c_{\min}+\alpha(c_{\max}-c_{\min})\}.
\]

Um candidato da RCL é selecionado aleatoriamente (uniforme), aplicado, e o processo continua até não existir inserção factível.

\paragraph{Compatibilidade com o baseline I1.}
Embora o GRASP use aleatorização via RCL, ele preserva a mesma estrutura de construção por inserção e os mesmos testes de factibilidade (capacidade e janelas) utilizados no restante do pipeline.
Isso evita que o ganho observado seja resultado de regras de viabilidade diferentes.

\subsection{GRASP fixo}

No modo fixo, o valor de $\alpha$ permanece constante em toda a execução.
No código, esse valor é fornecido por \texttt{--alpha} (padrão 0.3).

\paragraph{Interpretação.}
Valores menores de $\alpha$ tornam a construção mais gulosa (intensificação);
valores maiores ampliam diversidade (exploração).
A calibração desse parâmetro pode alterar significativamente o balanço qualidade/tempo.

\subsection{GRASP reativo}

No modo reativo, $\alpha$ é amostrado de um conjunto discreto e suas probabilidades são atualizadas periodicamente \cite{prais2000reactive}.
Seja $\bar{f}_i$ o custo médio observado para o valor $\alpha_i$ no bloco atual e $f^*$ o melhor custo global até o momento.
A qualidade relativa é calculada por:
\[
q_i = \left(\frac{f^*}{\bar{f}_i}\right)^\gamma,
\qquad
p_i = \frac{q_i}{\sum_j q_j}.
\]

O parâmetro $\gamma$ controla quão agressiva é a concentração de probabilidade nos $\alpha$ de melhor desempenho.

\subsection{Parâmetros relevantes}

Os parâmetros principais da implementação são:
\begin{itemize}
  \item comuns: $N_{\max}$ (\texttt{maxIters}) e $L$ (\texttt{maxNoImprove});
  \item modo fixo: $\alpha$;
  \item modo reativo: conjunto \texttt{reactiveAlphas}, \texttt{reactiveUpdate} e $\gamma$.
\end{itemize}

Esses parâmetros compõem o espaço de calibração no fluxo com \emph{irace}, mantendo o construtivo base (I1) fixo para comparação justa entre algoritmos.
